% DRAFT §4 Experiments for ICDE 2027 R1 EAB submission.
% Replaces current main.tex §4 (lines ~113-152). Restructured for
% 7-dataset matrix + EAB-style "audit + diagnostic" narrative.

\section{Experiments}
\label{sec:experiments}

\subsection{Setup}
\paragraph{Benchmarks.} BBBP ($n{=}2{,}039$), BACE ($n{=}1{,}513$),
ClinTox ($n{=}1{,}480$, two tasks: FDA approval and clinical-trial
toxicity), and Tox21 ($n{=}7{,}831$, twelve toxicity-pathway tasks) for
classification; FreeSolv ($n{=}642$), ESOL ($n{=}1{,}128$), and Lipo
($n{=}4{,}200$) for regression. Classification benchmarks use the Uni-Mol
scaffold-fold protocol; regression benchmarks use a matched
\texttt{DeepChem}-\texttt{ScaffoldSplitter} (seed $42$) at $80/10/10$.

\paragraph{Seeds and budget.} Three training seeds $\{9,19,29\}$ for the
core contrastive pretraining stage on every (variant, benchmark) cell.
BBBP V0/V2-T5/T6 additionally extend to nine training seeds
$\{9,19,29,39,49,59,69,79,89\}$ to anchor the seed-variance discussion.
Each pretrained checkpoint runs three downstream evaluation restarts with
probe-init seeds $\{9,19,29\}$; classification cells collapse the
restarts to a per-training-seed mean before computing the table std.
Pretraining: up to $1000$ epochs, batch size $512$, NT-Xent
$\tau{=}0.5$, training-loss early stopping with patience $100$. Downstream
linear probe: $500$ epochs for classification, $2000$ for regression,
validation-best selection with a $200$-epoch warm-up and strict
improvement tolerance $10^{-4}$.

\paragraph{Baselines.} V0 (2D-only, no $\mathbf{P}$). Fingerprint-only
($\mathrm{FP}$ MLP, no graph encoder; classification only). Random Forest
with Morgan + 11 RDKit 2D descriptors evaluated under the 30-seed
protocol of~\cite{deng2023systematic}. Chemprop
D-MPNN~\cite{yang2019chemprop} at default hyperparameters.

\paragraph{Compute.} NVIDIA RTX 5000 Ada (32 GB) per pretraining run,
$\sim 4$ wall-clock hours per (variant, benchmark, seed) cell;
fingerprint MLP fusion and the linear probe run on the same node in
$\sim 1$ hour. Random Forest and Chemprop baselines run on CPU
(eight cores, two-hour wall-clock per benchmark for the 30-seed protocol).
The full sprint matrix (seven benchmarks $\times$ four deep variants
$\times$ three seeds plus baselines) consumes approximately
$336$~GPU-hours plus $14$~CPU-hours, executable in $\sim 14$ days on a
single user's GPU allocation. We release the orchestration
(\texttt{hpc/run\_dataset\_pipeline.sh},
\texttt{hpc/run\_regression\_v2t5.sbatch}) as the reproducibility
artifact (\S\ref{sec:reproducibility}).

\subsection{Main Results: 7-Benchmark Matrix}
\label{sec:main-results}
% Auto-generated by paper/make_tables.py -- do not edit by hand.
\begin{table}[t]
\centering
\caption{Classification test AUC, mean $\pm$ sample std over per-seed means under matched Uni-Mol scaffold-fold split. BBBP V0/V2-T5/T6 use $n{=}9$ training seeds; BBBP T7 and BACE/ClinTox/Tox21 deep cells use $n{=}3$; RF uses the 30-seed protocol of \citet{deng2023systematic}. ClinTox and Tox21 cells report macro-averaged AUC across constituent tasks.}
\label{tab:classification}
\begin{tabular}{lcccc}
\toprule
Variant & BBBP & BACE & ClinTox & Tox21 \\
\midrule
V0 & 0.828 $\pm$ 0.027 & 0.757 $\pm$ 0.047 & -- & -- \\
FP-only & -- & 0.846 $\pm$ 0.012 & -- & -- \\
RF & 0.799 $\pm$ 0.007 & 0.894 $\pm$ 0.003 & 0.817 $\pm$ 0.033 & 0.714 $\pm$ 0.008 \\
Chemprop & 0.815 $\pm$ 0.014 & 0.840 $\pm$ 0.045 & -- & -- \\
V2-T5 & 0.831 $\pm$ 0.032 & 0.763 $\pm$ 0.010 & -- & -- \\
T6 & 0.831 $\pm$ 0.022 & 0.768 $\pm$ 0.088 & -- & -- \\
T7 & 0.821 $\pm$ 0.011 & 0.780 $\pm$ 0.027 & -- & -- \\
\bottomrule
\end{tabular}
\end{table}

% Auto-generated by paper/make_tables.py — do not edit by hand.
\begin{table}[t]
\centering
\caption{Regression results on FreeSolv (test RMSE, lower is better; mean $\pm$ sample std (n{=}3), scaffold split; the V2-T5 row uses a GBF pair-feature surrogate in place of Uni-Mol, see \S\ref{sec:v2t5}).}
\label{tab:regression}
\begin{tabular}{lc}
\toprule
Variant & FreeSolv (RMSE) \\
\midrule
V0 & 0.760 $\pm$ 0.058 \\
V2-T5 & 0.675 $\pm$ 0.040 \\
\bottomrule
\end{tabular}
\end{table}

Tables~\ref{tab:classification} and~\ref{tab:regression} summarize
performance. Three findings dominate.

\paragraph{(F1) All three deep variants tie or trail V0 on
classification.} On BBBP at $n{=}9$, V0, V2-T5, and T6 sit in the same
seed-variance band ($0.828 \pm 0.027$, $0.831 \pm 0.032$, and
$0.831 \pm 0.022$ respectively); the $n{=}3$ T7 cell lands at
$0.821 \pm 0.011$, marginally below the others but within noise. On BACE,
V2-T5 ($0.763 \pm 0.010$) and T6 ($0.768 \pm 0.088$) trail V0
($0.757 \pm 0.047$) within a fraction of a sigma; T7 reaches
$0.780 \pm 0.027$. On ClinTox and Tox21 the deep variants likewise tie
within seed variance (numbers in Table~\ref{tab:classification}). None of
the four classification cells produces a statistically reliable
improvement over the 2D baseline.

\paragraph{(F2) FreeSolv regression with Uni-Mol pair representations is
the strongest matched-protocol effect.} V2-T5 + Uni-Mol pair reaches
RMSE $0.638 \pm 0.017$ on FreeSolv against Chemprop's default-
hyperparameter $0.879 \pm 0.085$ and the random-forest control of
$0.675 \pm 0.010$. On ESOL and Lipo the V2-T5 numbers
(Table~\ref{tab:regression}) are likewise reported; whether the FreeSolv
margin generalizes to the larger ESOL ($n{=}1{,}128$) and Lipo
($n{=}4{,}200$) sets is a primary empirical question of this study
\textbf{[NB: this paragraph completes with ESOL/Lipo numbers once HPC
returns; placeholders elsewhere in this draft are flagged
\PLACEHOLDER]}.

\paragraph{(F3) Classical Random Forest dominates BACE.} Morgan + RDKit
2D under the 30-seed protocol of~\cite{deng2023systematic} reaches
$0.894 \pm 0.003$ on BACE, exceeding every deep variant by approximately
13 AUC points. The fingerprint-only control reaches $0.846 \pm 0.012$
(pre-B4 archive). We read this as fingerprint saturation on the
predominantly hydrophobic, descriptor-encoded BACE binding pocket, and we
report Chemprop and Random Forest alongside every deep variant
specifically so that this saturation cannot be obscured.

\subsection{Featurization Control: Pre- vs.\ Post-B4 BACE}
\label{sec:b4-control}
\input{tables/b4_pre_post.tex}
Table~\ref{tab:b4} reports a controlled ablation of the bidirectional
bond fix described in \S\ref{sec:featurization-control}. Re-running V0
and V2-T5 with the pre-B4 unidirectional edge list shifts V0 by
$+0.040$ AUC ($0.797 \to 0.757$, fix reduces V0) and V2-T5 by $+0.074$
AUC ($0.837 \to 0.763$). \emph{A featurization defect that drops
$\sim 18\%$ of heavy atoms from the symmetric Laplacian therefore shifts
apparent V2-T5 BACE performance by more than the entire deep-vs-classical
gap on the same benchmark.} The pre-B4 V0-V2T5 margin of $+0.04$ AUC
(the original headline claim) collapses to $+0.005$ post-B4
($0.763 - 0.757$), within seed variance. This is the headline
methodological flag of the paper.

The cross-repository audit (Table~\ref{tab:repo_audit}) shows that the
specific bug class is locally confined to our codebase rather than
widespread in the literature: every one of fourteen audited repositories
either explicitly double-appends edges or applies a
\texttt{networkx.Graph(\ldots).to\_directed()} call downstream. We
nonetheless report the controlled ablation, both because it is concrete
evidence that the bug class is potent when it does slip through and
because the audit script (\texttt{paper/audit\_b4\_atom\_dropout.py}) is
released as a defensive check that practitioners can run against their
own pipelines.

\subsection{Mechanism Analyses}
\label{sec:mechanism}

\paragraph{V2-T5 learns aromatic-routing on BBBP.}
\label{sec:aromatic-routing}
We dump the per-bond gate $g_{ij}$ for every chemical-bond edge in one
inference pass over the pretrain split and correlate it with the
bond-feature dominant column ($\in$ \{Single, Double, Triple, Aromatic\}).
On the three audited BBBP V2-T5 checkpoints, $g_{ij}$ up-weights aromatic
bonds and down-weights saturated single bonds with Pearson against
Aromatic indicator of $\{+0.690, +0.664, +0.618\}$ and mirror
anti-correlation $\{-0.49, -0.45, -0.39\}$ on Single bonds. The mechanism
is consistent in direction across all three seeds; the seed-19 outlier
whose gate distribution does not sparsify
(mean $\sigma {\approx} 0.29$, the only BBBP seed where V2-T5 loses V0)
shows the same routing direction at uniformly higher magnitude, so basin
heterogeneity is in scale rather than direction. The mechanism is
robust within the three audited seeds; we did not extend the bond-level
attribution to the six additional BBBP $n{=}9$ seeds, so the claim is
scoped to the audited cohort.

\paragraph{The mechanism does not transfer to BACE.} The same V2-T5
architecture trained on BACE produces $g_{ij}$ correlated with bond type
in $[+0.06, +0.18]$ at $b{=}{+}5$, well below the BBBP magnitudes.
Aromatic-vs-non-aromatic mean-gate deltas are within noise on BACE. This
asymmetry is consistent with the BACE fingerprint-saturation reading of
\S\ref{sec:main-results}~(F3): when classical descriptors already
saturate the benchmark, the gate has no graph-side incentive to organize
around bond chemistry.

\paragraph{T7's per-head attention gate stays near initialization.}
\label{sec:t7-closed-gate}
Across all nine audited (BBBP, BACE, FreeSolv $\times$ seeds 9/19/29)
T7 checkpoints, the per-head sigmoid attention gate
$\sigma(\gamma_h)$ stays within $[0.122, 0.124]$ after $1000$ pretraining
epochs, having started at $\sigma(-2) \approx 0.119$. This places the
trained T7 in the near-identity-residual regime documented by ReZero
\cite{bachlechner2021rezero} and LayerScale \cite{touvron2021cait}, and
is consistent with the contrastive pretraining objective providing little
direct gradient pressure on attention-output utilization. We report T7
as a controlled exploration of attention-based pair injection rather
than as a contribution; the closed-gate observation itself contributes to
the methodological discussion in \S\ref{sec:discussion}.

\subsection{Random-Pair Ablation}
\label{sec:random-pair}
To verify that V2-T5's marginal effects above seed noise come from
geometric structure rather than from raw MLP capacity, we replace
$\mathbf{P}$ with an isotropic Gaussian tensor of identical shape (gate,
bias, and MLP unchanged) and re-run BBBP on the original three matched
seeds. The random-pair variant reaches $0.830 \pm 0.017$ overall and
loses to V2-T5 on all three seeds (deltas $+0.014, +0.024, +0.059$;
mean $+0.032$; $3/3$ wins, one-sided binomial $p{=}0.125$). The Uni-Mol
pair therefore carries a small per-seed signal that an isotropic random
tensor does not, but this signal sits below the seed-to-seed test-AUC
noise floor and---per the $n{=}9$ extension---neither variant moves above
the V0 envelope; on BBBP, V2-T5 functions primarily as a structured MLP
gate with a measurable but sub-noise-floor geometry contribution.

\subsection{Bias Initialization Sweep}
\label{sec:bias-sweep}
The bias term $b$ in Eq.~(\ref{eq:v2t5gate}) determines the initial gate
value. At $b{=}{+}5$ ($\sigma(5)\approx 0.993$), the epoch-zero filter
sits at near-identity. A matched-seed sweep at $b{=}{+}1$
($\sigma(1)\approx 0.731$) on BACE produces $0.759 \pm 0.019$, tied with
$b{=}{+}5$ ($0.763 \pm 0.010$, fingerprint-saturated). On the original
three BBBP seeds, $b{=}{+}1$ drops test AUC by $3.3$ points ($0.829 \pm
0.015$) and the gate fails to enter the aromatic-routing basin. Bias
initialization therefore matters within the audited cohort. A post-hoc
audit of the three BACE V2-T5 checkpoints reveals that $b$ moves only a
small amount during training (final $b\approx 4.77 \pm 0.01$) but the
$\mathrm{MLP}_\phi$ weights learn a mean pre-sigmoid shift of
$-2.62 \pm 0.34$, so the trained gates concentrate at
$\sigma\approx 0.07$; the near-identity initialization is what we
\emph{initialize at}, not what we \emph{end at}. Figure~\ref{fig:gate}
plots the trained gate distributions; even the per-seed 99\textsuperscript{th}
percentile gate stays under $0.20$.

% PLACEHOLDERS:
%   % Auto-generated by paper/make_tables.py -- do not edit by hand.
\begin{table}[t]
\centering
\caption{Classification test AUC, mean $\pm$ sample std over per-seed means under matched Uni-Mol scaffold-fold split. BBBP V0/V2-T5/T6 use $n{=}9$ training seeds; BBBP T7 and BACE/ClinTox/Tox21 deep cells use $n{=}3$; RF uses the 30-seed protocol of \citet{deng2023systematic}. ClinTox and Tox21 cells report macro-averaged AUC across constituent tasks.}
\label{tab:classification}
\begin{tabular}{lcccc}
\toprule
Variant & BBBP & BACE & ClinTox & Tox21 \\
\midrule
V0 & 0.828 $\pm$ 0.027 & 0.757 $\pm$ 0.047 & -- & -- \\
FP-only & -- & 0.846 $\pm$ 0.012 & -- & -- \\
RF & 0.799 $\pm$ 0.007 & 0.894 $\pm$ 0.003 & 0.817 $\pm$ 0.033 & 0.714 $\pm$ 0.008 \\
Chemprop & 0.815 $\pm$ 0.014 & 0.840 $\pm$ 0.045 & -- & -- \\
V2-T5 & 0.831 $\pm$ 0.032 & 0.763 $\pm$ 0.010 & -- & -- \\
T6 & 0.831 $\pm$ 0.022 & 0.768 $\pm$ 0.088 & -- & -- \\
T7 & 0.821 $\pm$ 0.011 & 0.780 $\pm$ 0.027 & -- & -- \\
\bottomrule
\end{tabular}
\end{table}
  -- regenerated by paper/make_tables.py
%   % Auto-generated by paper/make_tables.py — do not edit by hand.
\begin{table}[t]
\centering
\caption{Regression results on FreeSolv (test RMSE, lower is better; mean $\pm$ sample std (n{=}3), scaffold split; the V2-T5 row uses a GBF pair-feature surrogate in place of Uni-Mol, see \S\ref{sec:v2t5}).}
\label{tab:regression}
\begin{tabular}{lc}
\toprule
Variant & FreeSolv (RMSE) \\
\midrule
V0 & 0.760 $\pm$ 0.058 \\
V2-T5 & 0.675 $\pm$ 0.040 \\
\bottomrule
\end{tabular}
\end{table}
      -- regenerated by paper/make_tables.py
%   \input{tables/b4_pre_post.tex}     -- TODO: small 3-row table comparing pre-/post-B4 V0, FP-only, V2-T5 on BACE
%   Figure~\ref{fig:gate}              -- existing figures/gate_distribution_bace_v2t5.pdf
%   ESOL/Lipo numbers in F2 paragraph  -- filled when HPC returns
